\chapter{Cái tôi và sai lầm}
\markboth{Cái tôi và sai lầm}{Ego and Mistakes}

\section{Người đàn ông mất việc nhưng vẫn không chịu học lại từ đầu.}
\begin{truyen}{Người đàn ông mất việc nhưng vẫn không chịu học lại từ đầu.}{Ego delays recovery.}
\chuhoa{A}{nh nghĩ học lại kỹ năng mới là việc dành cho người trẻ hơn và vị trí thấp hơn. Càng giữ thể diện, anh càng chậm thích nghi. Đến khi chấp nhận bắt đầu lại, anh mới thấy cái tôi đã lấy của mình nhiều thời gian hơn chính thất bại ban đầu.}
\textit{(He thought learning new skills was for younger or lower-ranked people. The more he protected his pride, the slower he adapted. When he finally started again, he saw that ego had stolen more time than the original failure.)}
\end{truyen}
\begin{baihoc}
Sai lầm không giữ ta lại lâu bằng lòng tự ái không chịu sửa sai.
\textit{(Mistakes do not hold us back as long as pride that refuses correction.)}
\end{baihoc}
\begin{ghinhoanh}
\textbf{Short English to remember:}

Ego delays recovery.

Start again before pride hardens.
\end{ghinhoanh}
\begin{tuhocnhanh}
1. Em đang chậm sửa điều gì chỉ vì tự ái? \textit{(What are you slow to fix because of pride?)}

2. Nếu bỏ bớt cái tôi, bước sửa đầu tiên của em là gì? \textit{(If you lowered your ego, what would be your first repair step?)}
\end{tuhocnhanh}
\ngancach

\section{Một người không xin lỗi vì sợ mình thua.}
\begin{truyen}{Một người không xin lỗi vì sợ mình thua.}{Pride can ruin what apology could save.}
\chuhoa{C}{ô biết mình đã nói quá lời, nhưng nghĩ xin lỗi sẽ làm mình yếu thế. Đến khi mối quan hệ lạnh hẳn, cô mới hiểu có những cái thua rất nhỏ cứu được những mất mát rất lớn.}
\textit{(She knew her words had gone too far, but thought apologizing would make her weak. When the relationship went cold, she understood that some small losses can prevent very great losses.)}
\end{truyen}
\begin{baihoc}
Nhiều người thua không phải vì xin lỗi, mà vì chần chừ xin lỗi quá lâu.
\textit{(Many people do not lose because they apologize, but because they wait too long to do it.)}
\end{baihoc}
\begin{ghinhoanh}
\textbf{Short English to remember:}

Pride can ruin what apology could save.

Say sorry before distance grows.
\end{ghinhoanh}
\begin{tuhocnhanh}
1. Em có đang giữ một lời xin lỗi chỉ vì không muốn hạ mình không? \textit{(Are you withholding an apology because you do not want to lower yourself?)}

2. Mối quan hệ nào đáng để em bỏ cái tôi xuống trước? \textit{(Which relationship is worth lowering your ego for?)}
\end{tuhocnhanh}
\ngancach

\section{Người trẻ nghĩ mình luôn đúng vì đọc được nhiều.}
\begin{truyen}{Người trẻ nghĩ mình luôn đúng vì đọc được nhiều.}{Knowledge is not humility.}
\chuhoa{C}{ậu có lý lẽ sắc, có bài viết hay trích dẫn nhiều. Nhưng đi lâu trong đời, cậu mới nhận ra hiểu biết mà thiếu khiêm tốn chỉ khiến người ta nhanh phán xét hơn, chứ chưa chắc sáng suốt hơn.}
\textit{(He had sharp arguments and many quotations. But after living longer, he realized that knowledge without humility often makes people quicker to judge, not necessarily wiser.)}
\end{truyen}
\begin{baihoc}
Biết nhiều mà không biết mình còn giới hạn thì vẫn là một kiểu non nớt.
\textit{(Knowing much without knowing your limits is still a form of immaturity.)}
\end{baihoc}
\begin{ghinhoanh}
\textbf{Short English to remember:}

Knowledge is not humility.

Read more, judge slower.
\end{ghinhoanh}
\begin{tuhocnhanh}
1. Em có đang dùng hiểu biết để hiểu người hơn hay để thắng người hơn? \textit{(Are you using knowledge to understand people better or just to defeat them?)}

2. Điều gì giúp em giữ được sự khiêm tốn? \textit{(What helps you stay humble?)}
\end{tuhocnhanh}
\ngancach

\section{Một người cố giữ hình ảnh giỏi giang nên không dám hỏi.}
\begin{truyen}{Một người cố giữ hình ảnh giỏi giang nên không dám hỏi.}{Image can block growth.}
\chuhoa{A}{nh sợ hỏi sẽ bị nghĩ là kém. Vì thế anh im lặng trong nhiều điều không rõ, và tiến chậm hơn người chịu học công khai. Sau này anh mới hiểu giữ hình ảnh quá kỹ thường làm ta chậm lớn thật sự.}
\textit{(He feared that asking would make him look weak. So he stayed silent through many unclear things and grew more slowly than those willing to learn openly. Later he understood that protecting image too carefully often slows real growth.)}
\end{truyen}
\begin{baihoc}
Ngại ngu trước mặt người khác thường khiến ta ngu lâu hơn cần thiết.
\textit{(Being afraid to look foolish often keeps us foolish longer than necessary.)}
\end{baihoc}
\begin{ghinhoanh}
\textbf{Short English to remember:}

Image can block growth.

Ask now. Learn faster.
\end{ghinhoanh}
\begin{tuhocnhanh}
1. Em đang im lặng ở đâu chỉ để giữ hình ảnh? \textit{(Where are you staying silent just to protect your image?)}

2. Câu hỏi nào em nên hỏi ngay? \textit{(What question should you ask right away?)}
\end{tuhocnhanh}
\ngancach

\section{Người ta chỉ thấy cái tôi của mình sau khi làm đau người khác.}
\begin{truyen}{Người ta chỉ thấy cái tôi của mình sau khi làm đau người khác.}{Ego is loudest in hindsight.}
\chuhoa{N}{hiều năm sau một cuộc đổ vỡ, anh xem lại cách mình nói, cách mình coi thường cảm xúc người khác, cách mình luôn đòi được hiểu mà ít khi chịu hiểu. Lúc đó anh mới nhận ra cái tôi ngày trước đã nói to đến mức nuốt hết tình thương.}
\textit{(Years after a broken relationship, he looked back at how he spoke, how he dismissed others' feelings, and how he always demanded understanding without giving it. Then he realized his ego had once spoken so loudly that it swallowed tenderness.)}
\end{truyen}
\begin{baihoc}
Cái tôi càng lớn, khả năng lắng nghe càng nhỏ; đến khi hiểu ra thì thường đã làm ai đó tổn thương.
\textit{(The bigger the ego, the smaller the ability to listen; by the time we realize it, someone has often already been hurt.)}
\end{baihoc}
\begin{ghinhoanh}
\textbf{Short English to remember:}

Ego is loudest in hindsight.

Listen before regret teaches you.
\end{ghinhoanh}
\begin{tuhocnhanh}
1. Em có đang đòi được hiểu nhiều hơn đang cố hiểu người khác không? \textit{(Are you demanding understanding more than trying to understand others?)}

2. Em cần lắng nghe ai bằng thái độ khác đi? \textit{(Whom do you need to listen to differently?)}
\end{tuhocnhanh}
\ngancach